% engineering/constraints_and_rules.tex

\chapter{Constraint Code Architecture}

\section*{Motivation}

\begin{remark}
Railway alignment rules must be represented in a form that is both
mathematically meaningful and computationally tractable.
\end{remark}

\begin{remark}
Within the AIM framework, this is achieved by expressing all rules as
constraint residuals that can be evaluated uniformly inside an
optimization loop.
\end{remark}

\begin{remark}
The implementation goal is therefore not a collection of isolated rule
checks, but a modular constraint library.
\end{remark}

\section{General Principle}

\begin{definition}[Constraint residual]
A constraint is represented as a function
\[
c(\mathbf{x}) \in \R^m
\]
whose value vanishes in the ideal case and whose magnitude measures
violation.
\end{definition}

\begin{remark}
This unified representation allows the same mathematical object to be
used for:
\begin{itemize}
\item feasibility checks,
\item penalty terms,
\item direct inclusion into SQP solvers.
\end{itemize}
\end{remark}

\begin{verbatim}
function evaluateConstraint(ctx) {
  return [residual1, residual2, ...]
}
\end{verbatim}

\section{Constraint Registry}

\begin{remark}
A practical implementation should organize constraints in a registry.
\end{remark}

\begin{verbatim}
const ConstraintRegistry = {
  geometric: {},
  dynamic: {},
  topologic: {},
  regulatory: {}
}
\end{verbatim}

\begin{remark}
This keeps domain logic extensible and avoids hard-coded monolithic
solver structures.
\end{remark}

\section{Constraint Object Model}

\begin{verbatim}
const constraint = {
  id: "min_radius",
  group: "regulatory",
  appliesTo: "edge",
  weight: 1.0,
  evaluate: (ctx) => []
}
\end{verbatim}

\begin{remark}
The fields have the following meaning:
\begin{itemize}
\item \texttt{id}: unique name,
\item \texttt{group}: semantic class,
\item \texttt{appliesTo}: edge, node, pair, network,
\item \texttt{weight}: optional scaling,
\item \texttt{evaluate}: residual callback.
\end{itemize}
\end{remark}

\section{Context Object}

\begin{remark}
Constraints should not operate directly on raw global arrays.
Instead, they receive a normalized context object.
\end{remark}

\begin{verbatim}
const ctx = {
  edge,
  node,
  pair,
  i,
  kappa,
  phi,
  geometry,
  speed,
  params,
  metadata
}
\end{verbatim}

\begin{remark}
This provides a stable interface between model, solver, and rule system.
\end{remark}

\section{Geometric Constraints}

\subsection{Position Continuity}

\begin{verbatim}
function continuityConstraint(p1, p2) {
  return [
    p1.x - p2.x,
    p1.y - p2.y
  ]
}
\end{verbatim}

\begin{remark}
This enforces \(C^0\)-continuity between adjacent geometric elements or
connected alignments.
\end{remark}

\subsection{Direction Continuity}

\begin{verbatim}
function directionConstraint(t1, t2) {
  return [
    t1.dx - t2.dx,
    t1.dy - t2.dy
  ]
}
\end{verbatim}

\begin{remark}
This corresponds to tangent continuity and therefore to \(C^1\)-style
behaviour in geometric terms.
\end{remark}

\subsection{Curvature Continuity}

\begin{verbatim}
function curvatureConstraint(k1, k2) {
  return [k1 - k2]
}
\end{verbatim}

\begin{remark}
This enforces continuity of curvature across boundaries.
\end{remark}

\section{Dynamic Constraints}

\subsection{Equilibrium Constraint}

\begin{verbatim}
function equilibriumConstraint(kappa, phi, v, g) {
  return [v * v * kappa - g * Math.sin(phi)]
}
\end{verbatim}

\begin{remark}
This residual measures cant deficiency or excess in analytic form.
\end{remark}

\subsection{Jerk Constraint}

\begin{verbatim}
function jerkConstraint(kappa0, kappa1, phi0, phi1, ds, v, g) {
  const dk = (kappa1 - kappa0) / ds
  const dphi = (phi1 - phi0) / ds
  const j = v * v * v * dk - g * v * Math.cos(phi0) * dphi
  return [j]
}
\end{verbatim}

\begin{remark}
This constraint reflects dynamic smoothness and comfort-related limits.
\end{remark}

\section{Regulatory Constraints}

\subsection{Minimum Radius}

\begin{verbatim}
function minRadiusConstraint(kappa, Rmin) {
  return [Math.abs(kappa) - 1 / Rmin]
}
\end{verbatim}

\begin{remark}
A non-positive value indicates compliance with the radius limit.
\end{remark}

\subsection{Maximum Cant}

\begin{verbatim}
function maxCantConstraint(phi, phiMax) {
  return [Math.abs(phi) - phiMax]
}
\end{verbatim}

\subsection{Cant Gradient}

\begin{verbatim}
function cantGradientConstraint(phi0, phi1, ds, limit) {
  return [Math.abs((phi1 - phi0) / ds) - limit]
}
\end{verbatim}

\section{Topological Constraints}

\subsection{Switch Node Constraint}

\begin{verbatim}
function switchConstraint(main, branch) {
  return [
    main.x - branch.x,
    main.y - branch.y,
    main.dx - branch.dx,
    main.dy - branch.dy
  ]
}
\end{verbatim}

\begin{remark}
This expresses the basic coupling at branching nodes.
\end{remark}

\subsection{Parallel Track Spacing}

\begin{verbatim}
function parallelConstraint(p1, p2, d) {
  const dx = p1.x - p2.x
  const dy = p1.y - p2.y
  return [Math.sqrt(dx * dx + dy * dy) - d]
}
\end{verbatim}

\section{Constraint Assembly}

\begin{remark}
The full residual vector is assembled from all active constraints.
\end{remark}

\begin{verbatim}
function buildConstraintResiduals(ctx, constraints) {
  const r = []

  for (const c of constraints) {
    const ri = c.evaluate(ctx)
    const w = Math.sqrt(c.weight ?? 1)
    for (const v of ri) r.push(w * v)
  }

  return r
}
\end{verbatim}

\begin{remark}
This produces a single solver-compatible vector independent of the
semantic source of the constraints.
\end{remark}

\section{Factory Functions}

\begin{remark}
Constraints should be created through factories rather than by writing
anonymous closures throughout the code base.
\end{remark}

\begin{verbatim}
function makeMinRadiusConstraint(Rmin, weight = 1) {
  return {
    id: "min_radius",
    group: "regulatory",
    appliesTo: "edge",
    weight,
    evaluate: ({ kappa, i }) => [
      Math.abs(kappa[i]) - 1 / Rmin
    ]
  }
}
\end{verbatim}

\begin{verbatim}
function makeJerkConstraint(v, g, ds, weight = 1) {
  return {
    id: "jerk",
    group: "dynamic",
    appliesTo: "edge",
    weight,
    evaluate: ({ kappa, phi, i }) => {
      const dk = (kappa[i + 1] - kappa[i]) / ds
      const dphi = (phi[i + 1] - phi[i]) / ds
      return [v * v * v * dk - g * v * Math.cos(phi[i]) * dphi]
    }
  }
}
\end{verbatim}

\section{Constraint Presets}

\begin{remark}
In practical railway design, different projects require different
constraint configurations.
\end{remark}

\begin{verbatim}
const presetHighSpeed = [
  makeMinRadiusConstraint(4000, 10),
  makeJerkConstraint(v, g, ds, 8),
  makeMaxCantConstraint(phiMax, 6)
]
\end{verbatim}

\begin{verbatim}
const presetExistingLine = [
  makeMinRadiusConstraint(800, 3),
  makeJerkConstraint(v, g, ds, 2),
  makeParallelConstraint(trackDistance, 5)
]
\end{verbatim}

\section{Integration into the Solver}

\begin{verbatim}
function buildResidualVector(ctx, constraints, objectives) {
  return [
    ...buildConstraintResiduals(ctx, constraints),
    ...buildObjectiveResiduals(ctx, objectives)
  ]
}
\end{verbatim}

\begin{remark}
This shows the central design principle of the architecture:
constraints and objectives share the same residual-based interface.
\end{remark}

\section{Discussion}

\begin{remark}
The proposed architecture has several advantages:
\begin{itemize}
\item modularity,
\item extensibility,
\item compatibility with least-squares solvers,
\item transparent mapping between engineering rules and code.
\end{itemize}
\end{remark}

\begin{remark}
Instead of scattering rule checks throughout the application, the
constraint library concentrates domain logic in a single mathematical
layer.
\end{remark}

\section{Summary}

\begin{summary}
The constraint library translates railway rules, topology, and dynamic
requirements into a unified residual-based architecture.

This provides the technical bridge between theoretical modelling,
engineering regulations, and computational optimization.

As a result, the AIM framework becomes not only mathematically coherent,
but also directly implementable in software.
\end{summary}
